\section{Later rows control a stronger multiplicity comparison}

\subsection*{Finding M-011: a cut through the nested tail support}
Let \(m\geq2\), \(n>m\), and let \(\lambda\vdash nd\)
have \(r\) rows. For \(2\leq k\leq r+1\), put
\(\lambda_{r+1}=0\), \(t=\min\{m,\lambda_k\}\), and
\[
 C_{m,k}(\lambda)=
 \binom{m+k-2}{k-2}
 +\binom{r-k+1+t}{t}-1.
\]
Then
\begin{equation}\label{eq:M011}
 n\geq C_{m,k}(\lambda)
 \quad\Longrightarrow\quad
                  q^m_\lambda(d[n])\leq K_n(\lambda).
\end{equation}
The multiplicities are those of the actual reduced orbit
closures, with the internal-padding convention. If the ambient
multiplicity is zero, the conclusion is immediate; otherwise
\(r\leq d\). The one-row case also has no obstruction.

Two uniform consequences, at every outer degree, are
\begin{align}
 n>3m^4, q^m_\lambda>K_n(\lambda)
       &\quad\Longrightarrow\quad \lambda_6\geq3,
                                                    \label{eq:M011IP}\\
 n\geq m^{25},\ q^m_\lambda>K_n(\lambda)
       &\quad\Longrightarrow\quad \lambda_{27}\geq13.
                                                    \label{eq:M011BIP}
\end{align}
These constrain later rows even when the second row is larger
than \(n\). They therefore apply beyond the stabilization
criteria that require \(n\geq\lambda_2\).

\paragraph{Support lemma.}
Every highest-weight vector of weight \(\lambda\) in
\(\operatorname{Sym}^d\operatorname{Sym}^n\mathbb C^r\)
uses only coefficient generators indexed by \(\eta\) with
\[
        |\eta|=n,\qquad
        \sum_{i=j}^r\eta_i\leq\lambda_j
                       \quad(2\leq j\leq r).          \tag{M-011.1}
\]
Here ``uses'' refers to every generator in every coefficient
monomial of the polynomial, not only the sum of their weights.

To prove this, realize the symmetric power as the column-permutation
invariants in the balanced column-degree \((n,\ldots,n)\)
polynomials on an \(r\)-by-\(d\) matrix. By the complete
flag-minor spanning argument in F-014, its row highest-weight
space is spanned by products of flag minors whose heights are
the column lengths of \(\lambda\). Exactly \(\lambda_j\)
minor factors have height at least \(j\). In any fixed column,
a monomial of a single minor uses at most one variable; a minor
of height less than \(j\) uses none from rows \(j,\ldots,r\).
Thus each matrix monomial has total degree at most \(\lambda_j\)
in these rows of that column. Taking linear combinations cannot
introduce a monomial that violates the bound. Finally, the
averaging map from \(\operatorname{Sym}^d\operatorname{Sym}^n\)
sends a monomial basis to disjoint orbit sums of column monomials.
The condition holds in every column and is permutation invariant.
It therefore implies precisely (M-011.1) on each coefficient
generator. This also proves the support assertion for every
multiplicity copy, as in F-016.

\begin{proof}[Proof of the comparison]
If \(q^m_\lambda=0\), the conclusion is immediate. Otherwise
the permanent length bound gives \(r\leq m^2<n^2\), so
the coordinate subspace used below exists.
Take a highest-weight vector \(w\) of weight \(\lambda\) in
\(I(\Omega_n)_d\). It depends only on the first \(r\)
coordinates. A general restricted translate of the padded
permanent is \(L^{n-m}f\), with \(f\) of degree \(m\).
On the dense open set \(L(e_1)\ne0\), the upper-unipotent
change of basis in M-009 aligns \(L\) with \(x_1\) and
preserves the vanishing of \(w\). Hence it is enough to
consider \(x_1^{n-m}f\).

Expand \(f\) in monomials
\(x_1^{m-|\alpha|}x_2^{\alpha_2}\cdots x_r^{\alpha_r}\).
Delete every term violating one of the inequalities
\(\sum_{i=j}^r\alpha_i\leq\lambda_j\), \(j\geq2\),
and denote the resulting form by \(f_{\mathrm{tail}}\).
Padding changes only the exponent of \(x_1\), so (M-011.1)
gives the exact equality
\[
 w(x_1^{n-m}f)=w(x_1^{n-m}f_{\mathrm{tail}}).
\]
Split the other variables into the prefix
\(x_2,\ldots,x_{k-1}\) and the suffix \(x_k,\ldots,x_r\).
All retained monomials have prefix degree at most \(m\)
and suffix degree at most \(t=\min(m,\lambda_k)\).

Use the variable-split graph of M-007 with all prefix
multi-indices of degree at most \(m\) and all suffix
multi-indices of degree at most \(t\). These sets are
downward closed. Identify the two roots, insert coefficient
edges only for monomials actually retained in
\(f_{\mathrm{tail}}\), and put the \(x_1\)-loops and
negative variable-edge labels as before. The number of vertices
is \(S=C_{m,k}(\lambda)\). Every nontrivial cycle has one
coefficient edge and the two unique tree paths. The cycle
permutation sign cancels the signs on those variable edges.
A monomial with tail degree \(|\alpha|\) contributes
\(c_\alpha x^\alpha x_1^{S-|\alpha|}\).
The relevant cycle has \(|\alpha|+1\) distinct vertices,
so this exponent is nonnegative, including when \(S<m\).
Adjoining \(x_1 I_{n-S}\) gives a determinant equal to
\(x_1^{n-m}f_{\mathrm{tail}}\) whenever \(n\geq S\).

The vector \(w\) vanishes on this determinant and hence
on the dense open subset of the permanent orbit just used.
It vanishes on the entire reduced orbit closure. This holds
for every highest-weight vector in the determinant ideal,
so the ideal multiplicity spaces satisfy
\(J_{\Omega_n,\lambda}\subseteq J_{Z_{n,m},\lambda}\).
Subtracting their dimensions from the common ambient
multiplicity proves \eqref{eq:M011}.
\end{proof}

\begin{proof}[Proof of the two uniform consequences]
A positive permanent multiplicity has \(r\leq m^2\).
In the IP range, M-008 already excludes \(r\leq9\), so
we may use \(k=6\). If \(\lambda_6\leq2\), then
\[
 C_{m,6}(\lambda)
 \leq\binom{m+4}{4}+\binom{r-3}{2}-1
 \leq m^4+r^2-1<3m^4.
\]
The middle bound uses \(m\geq2\): the ratio
\(\binom{m+4}{4}/m^4=\prod_{i=1}^4(1+i/m)/4!\)
decreases with \(m\) and is \(15/16\) at \(m=2\).
The comparison theorem therefore forbids an obstruction.

In the BIP range, M-008 already excludes \(r\leq51\),
so use \(k=27\). If \(\lambda_{27}\leq12\),
\[
 \begin{split}
 C_{m,27}(\lambda)
 &\leq\binom{m+25}{25}+\binom{r-14}{12}-1\\
 &\leq\tfrac12m^{25}+m^{24}-1<m^{25}.
 \end{split}
\]
For the first term,
\(\binom{m+25}{25}/m^{25}\) decreases with \(m\),
and its value at \(m=2\) is \(351/2^{25}<1/2\).
For the second term, the monomials of degree at most twelve
in \(r-26\) variables inject into words of length twelve
on an alphabet of \(r-25\leq r\) letters, after inserting
a slack variable. Their number is at most \(r^{12}\leq m^{24}\).
Finally \(m^{24}\leq m^{25}/2\). This proves
\eqref{eq:M011BIP}.
\end{proof}

\paragraph{A general power-range consequence.}
For every integer \(c\geq6\), the same proof gives
\[
 n\geq m^c,\ q^m_\lambda>K_n(\lambda)
 \quad\Longrightarrow\quad
       \lambda_{c+2}\geq\lfloor c/2\rfloor+1.          \tag{M-011.2}
\]
Indeed M-008 leaves at least \(2c+2\) rows, so take
\(k=c+2\), \(t_0=\lfloor c/2\rfloor\). If
\(\lambda_k\leq t_0\), the prefix count is at most
\(m^c/2\). This follows by the same decreasing-ratio
argument and \(\binom{c+2}{2}<2^{c-1}\) for \(c\geq6\);
the last inequality holds at six and remains true on
incrementing \(c\), since its left ratio is
\((c+3)/(c+1)<2\). The suffix count is at most
\[
 \binom{r-c-1+t_0}{t_0}
       \leq\frac{r^{t_0}}{t_0!}
       \leq\frac{m^{2t_0}}{t_0!}\leq\tfrac12m^c,
\]
since \(t_0\geq3\) and \(t_0\leq c+1\).
Their sum minus one is less than \(m^c\), proving the claim.

\paragraph{Grothendieck dictionary.}
The support lemma uses exterior-power minors of the universal
map between locally free sheaves, followed by the Cauchy
decomposition of its symmetric algebra; these are the Schur
operations of cote~133, scan~028, and SGA~6, VI,
Theorem~3.3, pp.~392--393. The suffix inequalities describe
finite torus-weight submodules, using SGA~3, I,
Proposition~4.7.3, pp.~35--36. Their affine coefficient
schemes map to the determinant locus by integral polynomial
identities. Highest-weight vanishing then gives an inclusion
of homogeneous ideal comodules on the characteristic-zero
fiber, in the representation/comodule dictionary of SGA~3,
I, Proposition~4.7.2, pp.~34--35. The resulting inequality
concerns coherent coordinate-ring quotients. It uses neither
an effective lift of an irreducible finite-group representation
to a set nor a saturation assumption.
